\section{主要工作}

本作品完成了从格密码原语、通用IB-ME协议到实际消息收发软件的完整系统。两种HIBE与五种常规
IBS形成十种组合，另接入LATTE＋Wang--PW目标规模组合。所有组合共用系统建立、双方密钥生成、
消息加密和匹配解密五个算法，并支持角色密钥保存、新进程恢复、命令行与浏览器操作。

\textbf{将目标规模环格采样组织为逐频率低秩计算。} Sageloli PARAM-I的相关扰动若展开为
$403456\times403456$双精度矩阵，单个矩阵的理论存储约为1.18~TiB。作品保留多项式表示，
在每个频率分解$65\times65$的Gram矩阵，完成PARAM-I、II的发钥、签名和公开验证；两档运行
实测峰值RSS分别约为1.36~GiB和6.19~GiB。具体算法和测量见第五章及第七章。

\textbf{用范围感知编码压缩原生IBS对象。} SRGG的公钥、机构状态、用户密钥和签名按字段取值范围
选择位宽，并采用规范二进制表示。三档参数中，机构状态长度降至对应JSON的18.3\%至19.8\%，
签名长度降至37.7\%至39.1\%；完整B组合采用紧凑签名格式。

\textbf{将固定身份对的委派移入接收会话。} Lai接收会话在准备阶段生成叶密钥，随后复用该密钥
处理连续消息。专项实验中，会话准备约510秒，后续解密中位时间由原路径的526秒降至28.6秒，
会话内解密提速约18.4倍、耗时减少约94.6\%。叶密钥编码长度为34.7~MiB。

十种常规组合在相同机器、源码、线程和输入下完成基本比较。表~\ref{tab:work-highlights}摘录
三组代表的在线收发指标；第七章给出全部组合、各阶段时间、对象长度和运行内存。

\begin{table}[H]
\centering\small
\caption{三组代表的同机在线收发结果}
\label{tab:work-highlights}
\begin{tabularx}{\textwidth}{@{}p{.30\textwidth}rrX@{}}
\toprule
\rowcolor{black!8} 代表组合 & 重复加密 & 重复解密 & 完整密文 \\
\midrule
A：LATTE＋ML-DSA & 13.5~ms & 59.4~ms & 139~KiB \\
B：LATTE＋SRGG & 11.3~s & 68.7~ms & 146~KiB \\
C：Lai＋Sageloli & 89.4~s & 409~s & 4.34~MiB \\
\bottomrule
\end{tabularx}
\begin{minipage}{.98\textwidth}\footnotesize
环境：Ultra~9~275HX，Sage~10.9，Python~3.13.14；双方身份各64~B，消息32~B。
每组1个独立系统，各项为固定密钥下第二次操作的实测值。C行采用逐次委派路径。
\end{minipage}
\end{table}

软件同时完成双方身份匹配、三类身份错配、密文字节篡改和对象恢复测试。十种组合共完成20次
诚实加密、30次诚实解密、10次新进程恢复和20次外层异常拒绝；网页入口支持消息发送、身份核对和解密结果查看。
